\section{Governance Recommendations}\label{sec:governance}
\answers{Addresses RQ6 as recommendations, not as an evaluated finding}

We distilled the findings into a checklist of thirty controls (GC-01 to GC-30). Each states its requirement; its rationale (threat classes, root causes, rubric dimensions, and whether it targets a vulnerability, a design weakness or a policy choice); how to verify it, with an exact \tool{} command or a manual check; and its mapping to NIST AI RMF~1.0~\cite{nistairmf}, NIST AI~600-1~\cite{nist600}, ISO/IEC~42001 Annex~A~\cite{iso42001}, OWASP~\cite{owaspllm2025,owaspagentic2026}, MITRE ATLAS~\cite{atlas} and the MAS AI risk-management guidelines~\cite{masairm2026}. The controls that follow most directly from our evidence are: treat repository configuration, instruction-file imports included, as untrusted code, and gate changes to it in CI (GC-08, GC-09, GC-23; RC2, E3 secret-in-state); lock approve-everything modes by managed policy, and regard allowlists as a convenience rather than a boundary (GC-06, GC-07; RC3); supply the sandbox and egress boundary yourself, since D02 and D03 score lowest (GC-16, GC-17, GC-20); never run a headless agent with secrets or write tokens on an untrusted trigger (GC-21, GC-22; E2 CI strata); keep secrets out of configuration files and the agent's environment, and restrict state directories (GC-14, GC-15; E5); and re-run the black-box probes and the canary bench on the exact version and model you deploy (GC-03, GC-04). The checklist has not been evaluated for coverage or usability; the OWASP MCP Top~10 remains in beta, the ISO/IEC~42001 Annex~A titles are taken from public listings, and NIST AI~600-1 is cited at document level where necessary.
